\chapter{Desarrollo y Diseño}
\label{sec:desarrollo}

\section{Función de Transferencia Requerida}

La transferencia buscada es una recta que pasa por los puntos $(V_{i1}; V_{o1}) = (\SI{2.73}{\volt}; \SI{0}{\volt})$
y $(V_{i2}; V_{o2}) = (\SI{3.23}{\volt}; \SI{5}{\volt})$:
\begin{gather*}
  V_o - V_{o1} = \frac{V_{o2} - V_{o1}}{V_{i2} - V_{i1}} \left( V_i - V_{i1} \right)
  \quad\Rightarrow\quad
  V_o = \frac{\SI{5}{\volt}}{\SI{0.5}{\volt}} \left( V_i - \SI{2.73}{\volt} \right)
\end{gather*}
\begin{equation}
  V_o = A_v \left( V_i - V_{REF} \right), \qquad A_v = 10, \qquad V_{REF} = \SI{2.73}{\volt}
  \label{eq:transferencia}
\end{equation}

Es decir, el CAS debe restar a la señal del sensor una tensión de referencia de $\SI{2.73}{\volt}$ y amplificar la
diferencia diez veces.

\section{Topología del Circuito Acondicionador}

Se eligió una configuración restadora ($U_3$ en la Figura~\ref{crkt:completo}), que resta primero la referencia y luego amplifica
la diferencia. La topología propuesta por la cátedra amplifica primero y luego resta, por lo que también se amplifica
la componente continua del sensor: la tensión intermedia supera los \SI{15}{\volt} ($10 \cdot \SI{2.73}{\volt}$), lo
cual no es posible con una alimentación simple de \SI{5}{\volt}.

Si se cumple la condición de apareamiento $R_3/R_4 = R_5/R_6$, la salida resulta:
\begin{equation}
  V_o = \frac{R_3}{R_4} \left( V_i - V_{REF} \right)
  \label{eq:restador}
\end{equation}
Comparando con la Ecuación~\ref{eq:transferencia}, $R_3/R_4 = A_v = 10$. Fijando $R_3 = R_5 = \SI{100}{\kilo\ohm}$ se
obtiene $R_4 = R_6 = \SI{10}{\kilo\ohm}$.

\begin{figure}[!ht]
  \centering
  \resizebox{0.72\textwidth}{!}{\begin{circuitikz}[american]
  % Amplificadores operacionales
  \node[op amp, noinv input up] (U1) at (3,4.5) {$U_1$};
  \node[op amp] (U3) at (10,4) {$U_3$};

  % Divisor de referencia y seguidor
  \draw (0,7.2) node[vcc]{$\SI{5}{\volt}$} to[R, l_={\shortstack{$R_1$\\ \SI{10}{\kilo\ohm}}}] (0,0 |- U1.+) coordinate (A)
    to[R, l_={\shortstack{$R_2$\\ \SI{12}{\kilo\ohm}}}] ++(0,-2) node[ground]{};
  \draw (A) to[short, *-] (U1.+);
  \draw (U1.-) -- ++(0,-0.8) -| (U1.out);

  % Etapa restadora
  \draw (U3.-) -- ++(-0.5,0) coordinate (m);
  \draw (U3.+) -- ++(-0.5,0) coordinate (p);
  \draw (U1.out) to[short, *-] ++(0.6,0) node[below]{$V_{REF}$} to[R, l={\shortstack{$R_4$ \SI{10}{\kilo\ohm}}}] ++(2.2,0) |- (m);
  \draw (m) to[short, *-] ++(0,1.4) coordinate (t) to[R, l={\shortstack{$R_3$ \SI{100}{\kilo\ohm}}}] (t -| U3.out) -- (U3.out);
  \draw (U3.out) to[short, *-o] ++(0.8,0) node[right]{$V_o$};

  % Polarizacion del sensor
  \coordinate (K) at (5.4,1.2);
  \draw (K) to[R, l={\shortstack{$R_7$\\ \SI{1.8}{\kilo\ohm}}}, *-] ++(0,1.6) node[vcc]{$\SI{5}{\volt}$};
  \draw (K) ++(0,-2) node[ground]{} to[zD, l=LM335] (K);
  \node[left] at (K) {$V_T$};
  \draw (K) to[R, l={\shortstack{$R_6$ \SI{10}{\kilo\ohm}}}] (K -| p) coordinate (q) -- (p);
  \draw (q) to[short, *-] ++(0,-0.3) to[R, l={\shortstack{$R_5$\\ \SI{100}{\kilo\ohm}}}] ++(0,-1.7) node[ground]{};
\end{circuitikz}
}
  \caption{Circuito acondicionador completo.}
  \label{crkt:completo}
\end{figure}

\section{Tensión de Referencia}

La referencia se obtiene de la alimentación de \SI{5}{\volt} mediante un divisor resistivo:
\begin{gather*}
  V_{REF} = \frac{R_2}{R_1 + R_2} \cdot \SI{5}{\volt}
  \quad\Rightarrow\quad
  R_1 = R_2 \left( \frac{\SI{5}{\volt}}{V_{REF}} - 1 \right) = 0.8315 \cdot R_2
\end{gather*}
La combinación de valores normalizados que mejor aproxima esta relación es $R_1 = \SI{10}{\kilo\ohm}$ y
$R_2 = \SI{12}{\kilo\ohm}$, con la que $V_{REF} = \SI{2.727}{\volt}$. Como la resistencia de Thévenin del divisor
($\SI{5.45}{\kilo\ohm}$) es comparable con $R_4$, se agrega un seguidor de tensión ($U_1$) que evita que la
etapa restadora cargue al divisor y altere tanto la referencia como la condición de apareamiento.

\section{Polarización del Sensor}

$R_7$ se calcula para que circule $I_1 = \SI{1}{\milli\ampere}$ por el sensor en el peor caso, a \SI{50}{\celsius}.
La corriente total $I_T$ que entrega $R_7$ se divide entre el sensor y la rama $R_6$-$R_5$ de entrada al restador:
\begin{gather*}
  I_2 = \frac{V_{\SI{50}{\celsius}}}{R_6 + R_5} = \frac{\SI{3.23}{\volt}}{\SI{110}{\kilo\ohm}}
    = \SI{29.36}{\micro\ampere}
  \qquad
  I_T = I_1 + I_2 = \SI{1.02936}{\milli\ampere}
  \\
  R_7 = \frac{\SI{5}{\volt} - V_{\SI{50}{\celsius}}}{I_T} = \frac{\SI{1.77}{\volt}}{\SI{1.02936}{\milli\ampere}}
    = \SI{1719.5}{\ohm} \;\Rightarrow\; R_7 = \SI{1.8}{\kilo\ohm}
\end{gather*}
Con el valor normalizado, la corriente por el sensor varía entre $\SI{1.24}{\milli\ampere}$ (\SI{0}{\celsius}) y
$\SI{0.95}{\milli\ampere}$ (\SI{50}{\celsius}), dentro del rango de operación de $\SI{400}{\micro\ampere}$ a
$\SI{5}{\milli\ampere}$ que indica la hoja de datos del LM335.

\section{Elección del Amplificador Operacional}

Se utiliza el MCP6001, que admite alimentación simple de \SI{1.8}{\volt} a \SI{6}{\volt} y tiene entradas y salida
\textit{rail-to-rail}. Esto permite alimentarlo directamente entre GND y \SI{5}{\volt} y que la salida alcance valores
cercanos a \SI{0}{\volt} y \SI{5}{\volt}, necesarios en los extremos del rango.
